\documentclass[10pt,a4paper]{amsart}
\usepackage{amsmath,amssymb,amsthm,graphicx}
% For commutative diagrams we use Xy-pic
\usepackage[all]{xy}

% We use 2cell for 2-commutative diagrams.
\xyoption{2cell}
% Theorem environments.
%
\theoremstyle{plain}
\newtheorem{theorem}[subsection]{Theorem}
\newtheorem{proposition}[subsection]{Proposition}
\newtheorem{lemma}[subsection]{Lemma}

\theoremstyle{definition}
\newtheorem{definition}[subsection]{Definition}
\newtheorem{example}[subsection]{Example}
\newtheorem{exercise}[subsection]{Exercise}
\newtheorem{situation}[subsection]{Situation}

\theoremstyle{remark}
\newtheorem{remark}[subsection]{Remark}
\newtheorem{remarks}[subsection]{Remarks}

\numberwithin{equation}{subsection}

% Macros
%
\def\lim{\mathop{\mathrm{lim}}\nolimits}
\def\colim{\mathop{\mathrm{colim}}\nolimits}
\def\Spec{\mathop{\mathrm{Spec}}}
\def\Hom{\mathop{\mathrm{Hom}}\nolimits}
\def\Ext{\mathop{\mathrm{Ext}}\nolimits}
\def\SheafHom{\mathop{\mathcal{H}\!\mathit{om}}\nolimits}
\def\SheafExt{\mathop{\mathcal{E}\!\mathit{xt}}\nolimits}
\def\Sch{\mathit{Sch}}
\def\Mor{\mathop{\mathrm{Mor}}\nolimits}
\def\Ob{\mathop{\mathrm{Ob}}\nolimits}
\def\Sh{\mathop{\mathit{Sh}}\nolimits}
\def\NL{\mathop{N\!L}\nolimits}
\def\CH{\mathop{\mathrm{CH}}\nolimits}
\def\proetale{{pro\text{-}\acute{e}tale}}
\def\etale{{\acute{e}tale}}
\def\QCoh{\mathit{QCoh}}
\def\Ker{\mathop{\mathrm{Ker}}}
\def\Im{\mathop{\mathrm{Im}}}
\def\Coker{\mathop{\mathrm{Coker}}}
\def\Coim{\mathop{\mathrm{Coim}}}

% Boxtimes
%
\DeclareMathSymbol{\boxtimes}{\mathbin}{AMSa}{"02}

%
% Macros for moduli stacks/spaces
%
\def\QCohstack{\mathcal{QC}\!\mathit{oh}}
\def\Cohstack{\mathcal{C}\!\mathit{oh}}
\def\Spacesstack{\mathcal{S}\!\mathit{paces}}
\def\Quotfunctor{\mathrm{Quot}}
\def\Hilbfunctor{\mathrm{Hilb}}
\def\Curvesstack{\mathcal{C}\!\mathit{urves}}
\def\Polarizedstack{\mathcal{P}\!\mathit{olarized}}
\def\Complexesstack{\mathcal{C}\!\mathit{omplexes}}
% \Pic is the operator that assigns to X its picard group, usage \Pic(X)
% \Picardstack_{X/B} denotes the Picard stack of X over B
% \Picardfunctor_{X/B} denotes the Picard functor of X over B
\def\Pic{\mathop{\mathrm{Pic}}\nolimits}
\def\Picardstack{\mathcal{P}\!\mathit{ic}}
\def\Picardfunctor{\mathrm{Pic}}
\def\Deformationcategory{\mathcal{D}\!\mathit{ef}}
\usepackage{cleveref}
\crefname{originallemma}{original lemma}{original lemmas}
\Crefname{originallemma}{Original lemma}{Original lemmas}
\makeatletter
\newlabel{derived-lemma-two-ss-complex-functor}{{015J}{source}{}{sourcepointer-1}{}}
\newlabel{derived-lemma-two-ss-complex-functor@cref}{{[originallemma][0][]015J}{[1][0][]source}{}{}{}}
\makeatother
\crefname{originalreference}{original reference}{original references}
\Crefname{originalreference}{Original reference}{Original references}
\makeatletter
\newlabel{lemma-base-change-Rf-star-colim}{{0EZT}{source}{}{sourcepointer-2}{}}
\newlabel{lemma-base-change-Rf-star-colim@cref}{{[originalreference][0][]0EZT}{[1][0][]source}{}{}{}}
\makeatother
\crefname{originalreference}{original reference}{original references}
\Crefname{originalreference}{Original reference}{Original references}
\makeatletter
\newlabel{theorem-smooth-base-change}{{0EYU}{source}{}{sourcepointer-3}{}}
\newlabel{theorem-smooth-base-change@cref}{{[originalreference][0][]0EYU}{[1][0][]source}{}{}{}}
\makeatother
\title{Smooth base change for a limit of smooth schemes and torsion coefficients invertible on the base}
\author{The Stacks Project Authors}
\date{stacks-project:etale-cohomology.tex:local-snapshot-20261011; GNU FDL 1.2 or later}
\begin{document}
\maketitle
\noindent\textit{Editorial note.} All delivered mathematical text is the author's own, copied from the pinned Stacks Project chapter file etale-cohomology.tex; no statement or proof step was rewritten, reordered or removed.\par\medskip
\noindent\textit{Editorial note.} The counted proof is the author's complete argument: the proof is written out in full: the two spectral sequences, the identification of their entries from the two base change inputs, and the resulting identification of the abutments.\par\medskip
\noindent\textit{Editorial note.} The label lemma-smooth-base-change-general names the counted statement, and the proof environment immediately following it in the source is delivered verbatim as its proof.\par\medskip
\noindent\textit{Editorial note.} The counted statement is carried by the author's own lemma environment in the chapter source.\par\medskip
\noindent\textit{Editorial note.} The statement's fibre square is the author's and is delivered verbatim; the three cited inputs are kept as external pointers with their canonical Stacks tags.\par\medskip
\noindent\textit{Allowed standard prerequisites (editorial).}\begin{itemize}
\item The smooth base change theorem. Same chapter, the smooth base change theorem for a smooth morphism with torsion coefficients invertible on the base. Supplies the identification g\textasciicircum{}\{-1\}R\textasciicircum{}p f\_* of a single sheaf with the derived pushforward along f'.
\item Base change along a limit of smooth schemes. Same chapter, the base change lemma for a directed inverse limit of schemes smooth over the base. Supplies the extension of that identification to the limit situation.
\item The spectral sequence of a complex computing the pushforward. Derived Categories, the lemma giving the two spectral sequences E\_2\textasciicircum{}\{p,q\} = R\textasciicircum{}p f\_*H\textasciicircum{}q(E) and the corresponding one for f' converging to the derived pushforwards. Supplies the comparison of the two spectral sequences that yields the statement for complexes.
\end{itemize}\medskip
\ifdefined\hypertarget\hypertarget{sourcepointer-1}{}\fi
\noindent\textit{Source reference.} Original lemma 015J: sources/stacks-chow/tags.txt, source lines 1510--1510. Prerequisite pointer invoked at proof line(s) 18003; it names a standard fact of a sibling chapter that the counted argument uses, and every step of the counted proof built on it is retained verbatim in the delivered body.\par\medskip
\ifdefined\hypertarget\hypertarget{sourcepointer-2}{}\fi
\noindent\textit{Source reference.} Original reference 0EZT: sources/stacks-etale-cohomology/pinned-original-source.tex, source lines 16410--16410. Prerequisite pointer invoked at proof line(s) 18006; it names a standard fact of the same chapter that the counted argument uses, and every step of the counted proof built on it is retained verbatim in the delivered body.\par\medskip
\ifdefined\hypertarget\hypertarget{sourcepointer-3}{}\fi
\noindent\textit{Source reference.} Original reference 0EYU: sources/stacks-etale-cohomology/pinned-original-source.tex, source lines 17446--17446. Prerequisite pointer invoked at proof line(s) 18005; it names a standard fact of the same chapter that the counted argument uses, and every step of the counted proof built on it is retained verbatim in the delivered body.\par\medskip
\begin{lemma}
\label{lemma-smooth-base-change-general}
Let $S$ be a scheme. Let $S' = \lim S_i$ be a directed inverse
limit of schemes $S_i$ smooth over $S$ with affine transition
morphisms. Let $f : X \to S$ be quasi-compact and quasi-separated
and form the fibre square
$$
\xymatrix{
X' \ar[d]_{f'} \ar[r]_{g'} & X \ar[d]^f \\
S' \ar[r]^g & S
}
$$
Then
$$
g^{-1}Rf_*E = R(f')_*(g')^{-1}E
$$
for any $E \in D^+(X_\etale)$ whose cohomology sheaves $H^q(E)$
have stalks which are torsion of orders invertible on $S$.
\end{lemma}

\begin{proof}
Consider the spectral sequences
$$
E_2^{p, q} = R^pf_*H^q(E)
\quad\text{and}\quad
{E'}_2^{p, q} = R^pf'_*H^q((g')^{-1}E) = R^pf'_*(g')^{-1}H^q(E)
$$
converging to $R^nf_*E$ and $R^nf'_*(g')^{-1}E$.
These spectral sequences are constructed in
Derived Categories, Lemma \ref{derived-lemma-two-ss-complex-functor}.
Combining the smooth base change theorem
(Theorem \ref{theorem-smooth-base-change})
with Lemma \ref{lemma-base-change-Rf-star-colim} we see that
$$
g^{-1}R^pf_*H^q(E) = R^p(f')_*(g')^{-1}H^q(E)
$$
Combining all of the above we get the lemma.
\end{proof}

\end{document}
